\documentclass[11pt]{article}

% =====================================================================
% Energy-Based Route Ranking: Learned Selection over an Agent's
% Tool Library. Academic format mirrors the Unbrowse paper trilogy.
% Every reported number is a reproduced witness; the honest-negatives
% section is load-bearing, not decoration.
% =====================================================================

\usepackage[margin=1in]{geometry}
\usepackage{hyperref}
\usepackage{amsmath}
\usepackage{amssymb}
\usepackage{enumitem}
\usepackage{booktabs}
\usepackage{xcolor}
\hypersetup{colorlinks=true, linkcolor=blue, urlcolor=blue, citecolor=blue}

\title{\textbf{Energy-Based Route Ranking}\\
\large Learned Selection over an Agent's Tool Library}

\author{Lewis Tham\\ Unbrowse AI \\ \texttt{lewis@unbrowse.ai}
  \and Cayden Chik\\ Unbrowse AI \\ \texttt{cayden@unbrowse.ai}}
\date{June 2026}

\begin{document}
\maketitle

\begin{abstract}
An autonomous agent accumulates a library of callable routes --- learned
first-party API endpoints, each answering some class of intent. The runtime
question is selection: given a user intent and a library of candidate routes,
which route should fire? We treat selection as an energy-based scoring problem.
A learned energy head assigns a scalar compatibility energy $E(\text{intent},
\text{route})$ to each candidate; the selector ranks the library by ascending
energy and fires the lowest-energy match. The head is trained from observed
\emph{(intent, route, outcome)} triples and embedded into the runtime so that
ranking happens at inference, with a content-addressed cache so a repeated
selection is a lookup rather than a recomputation. We report only reproduced
results. On discrete-structure tasks the method wins: route ranking reaches
\textbf{R@1 $=0.0488$}, \textbf{4.6$\times$} a keyword baseline; access-pattern
type classification reaches \textbf{0.71--0.81 accuracy}, \textbf{$\approx$2.1$\times$}
baseline across two seeds; a closed-loop learned ranker reaches a cold-cell
\textbf{AUC of 0.750} against $0.5$ chance; the runtime-embedded head scores
\textbf{0.83 warm / 0.64 cold}; and a content-addressed cache yields a
\textbf{92$\times$ mean speedup} on a reuse micro-benchmark. We are equally
explicit about where the method gives nothing: energy-reranking a free-form
language model's own generations is indistinguishable from random
(lift $+0.000$), and on a public 101-route retrieval set a plain keyword
baseline beats the energy ranker. The honest boundary is the contribution: a
learned energy plus a content-addressed cache wins on \emph{discrete structure}
--- route retrieval, type classification, cache reuse --- and contributes nothing
to free-form prose generation.
\end{abstract}

\section{Introduction}
Once an agent learns callable routes from real usage, it holds a growing tool
library: each route is a typed, replayable first-party endpoint that answers some
class of intent. The hard runtime problem is no longer discovery but
\emph{selection} --- given an incoming intent and dozens or hundreds of candidate
routes, pick the one that actually answers it. A naive keyword or embedding match
treats the library as a flat lexical index and ignores the structured signal in
how a route was used and what it returned.

We frame selection as energy-based learning~\cite{lecun}. A learned energy head
maps an \emph{(intent, route)} pair to a scalar \emph{energy} --- low energy means
the route is compatible with the intent, high energy means it is not. The selector
embeds the intent, scores every candidate route, ranks the library by ascending
energy, and fires the lowest-energy route. The head is trained discriminatively
from observed \emph{(intent, route, outcome)} triples: routes that succeeded for
an intent are pushed to low energy, routes that did not are pushed up. The trained
head is small enough to embed directly in the runtime, so ranking is an inference
pass, and a content-addressed cache stores each result under the hash of its
content so a repeated selection resolves without recomputing.

We make one claim and bound it carefully. The method is a \emph{discrete-structure
ranker}: it wins where candidates carry separable structure --- distinct routes,
distinct access-pattern types --- and it wins by a wide margin there. It is not a
generative model and adds no signal to free-form text. This paper reports both
sides as reproduced measurements, and the negative side is stated as plainly as the
positive. Our contributions:
\begin{enumerate}[leftmargin=1.6em]
  \item An energy-based selector over an agent's route/tool library, trained from
  observed outcome triples and embedded into the runtime for inference-time
  ranking.
  \item A reproduced evaluation on discrete-structure tasks --- route ranking,
  access-pattern type classification, a closed-loop learned ranker, a
  runtime-embedded head, and content-addressed cache reuse --- each with a runnable
  gate.
  \item An explicit negative-results section establishing the method's boundary:
  it does nothing for free-form generation, and a keyword baseline can beat it on a
  public retrieval set. The boundary is the result.
\end{enumerate}

\section{Related work}
\subsection{Energy-based and joint-embedding models}
Energy-based models score the compatibility of an input pair with a single scalar
and select by minimising that scalar, rather than normalising a probability over a
large output space~\cite{lecun}. Joint-embedding predictive architectures extend
the idea to representation learning: predict the latent of a target from the latent
of a context and penalise the gap in embedding space~\cite{ijepa}. We borrow the
scoring-and-ranking shape --- score each candidate, rank by energy, pick the lowest
--- and apply it to a concrete, finite selection problem (which route fires) rather
than to representation pre-training.

\subsection{Retrieval and learned ranking}
Dense retrieval embeds queries and documents into a shared space and ranks by
similarity~\cite{dpr}; learning-to-rank optimises an ordering objective over
candidates~\cite{ltr}. Our setting differs in the training signal: candidates are
\emph{routes}, and the supervision is the observed \emph{outcome} of firing a route
for an intent, not a relevance label. The energy head is a discriminative ranker
over that outcome signal.

\subsection{Tool selection for agents}
Tool-use work teaches a model \emph{when} to call a tool~\cite{toolformer} and to
emit correct calls across many declared APIs~\cite{gorilla}. These rank tools
inside the language model's own forward pass. We separate selection from
generation: a dedicated energy head ranks the route library, and the language model
is left to do what it is good at. The negative-results section below is precisely
the boundary between these two responsibilities.

\section{Method}
\subsection{The selection problem}
Let an intent $q$ arrive against a library $\mathcal{R} = \{r_1, \dots, r_n\}$ of
candidate routes. A learned energy head $E_\theta$ scores each pair:
\begin{equation}
  s_i \;=\; E_\theta(q, r_i), \label{eq:energy}
\end{equation}
and the selector returns the lowest-energy route,
\begin{equation}
  r^\star \;=\; \arg\min_{r_i \in \mathcal{R}} E_\theta(q, r_i), \label{eq:argmin}
\end{equation}
or, for a shortlist, the library sorted by ascending $s_i$. Low energy encodes
``this route answers this intent.''

\subsection{Training the energy head}
The head is trained from observed \emph{(intent, route, outcome)} triples logged at
runtime. For an intent, the route that successfully answered it is the positive; a
sampled set of other library routes are negatives. The objective pushes the
positive's energy below the negatives' by a margin, so that
Eq.~\eqref{eq:argmin} selects the route that historically worked. Because the
supervision is the real outcome of firing a route, the head improves as the agent
accumulates usage --- the ranker is closed-loop.

\subsection{Embedding the head in the runtime}
The trained head is small and is loaded directly into the live selection path, so
ranking the library is one inference pass per resolution. A content-addressed cache
sits in front of execution: each intermediate result is stored under the hash of
its content, so a second identical selection resolves to a cached lookup instead of
re-scoring and re-firing. Selection cost is therefore paid once per distinct intent
and amortised over every repeat.

\subsection{The self-similar shape}
The same score-rank-select shape recurs at more than one scale. Selecting
\emph{which} route fires is one energy ranking; classifying a route's
\emph{access-pattern type} (e.g.\ anchor, server-rendered, structured-query) is the
same energy ranking over a small label set. A single learned-energy mechanism
serves both, which is why both appear in the evaluation.

\section{Evaluation}
We report only results backed by a runnable witness that re-runs green. Each row of
Table~\ref{tab:wins} names its gate; the appendix lists the reproduce commands.
Both the energy path and each baseline are scored on the same held-out data.

\subsection{Discrete-structure results (wins)}
\begin{table}[h]
\centering
\begin{tabular}{llrl}
\toprule
Task & Metric & Energy & Baseline \\
\midrule
Route ranking & R@1 (1 true / 99 distractors) & \textbf{0.0488} & 0.0106 \\
Access-pattern type & accuracy (2 seeds) & \textbf{0.71--0.81} & 0.32--0.39 \\
Closed-loop ranker & cold-cell AUC & \textbf{0.750} & 0.500 \\
Runtime-embedded head & warm / cold score & \textbf{0.83 / 0.64} & --- \\
Content-addressed cache & reuse speedup & \textbf{92$\times$} & 1$\times$ \\
\bottomrule
\end{tabular}
\caption{Reproduced discrete-structure results. Each row re-runs green under its
gate (Appendix~\ref{app:repro}).}
\label{tab:wins}
\end{table}

\paragraph{Route ranking.} Ranking the true route against 99 distractors for a real
intent, the energy head reaches \textbf{R@1 $=0.0488$ versus $0.0106$ for a keyword
baseline --- a 4.6$\times$ improvement}. The absolute number is small because the
task is severe (1-in-100 with real, near-duplicate distractors); the relevant
quantity is the multiple over the lexical baseline, and the learned energy more than
quadruples it.

\paragraph{Access-pattern type classification.} Energy-ranking a route's
access-pattern over a 7-class label set reaches \textbf{accuracy $0.71$--$0.81$
across two seeds, $\approx$2.1$\times$ a $0.32$--$0.39$ baseline}. The structure
here is discrete and separable, which is exactly the regime where the energy
ranker is strong.

\paragraph{Closed-loop learned ranker.} The closed-loop variant --- where the head
trains on accumulated outcomes and then ranks unseen cells --- reaches a
\textbf{cold-cell AUC of $0.750$ against $0.5$ chance}, confirming that a learned
ranker ships and generalises to cold cells, not just memorises warm ones.

\paragraph{Runtime-embedded head.} The head embedded into the live runtime scores
\textbf{$0.83$ warm and $0.64$ cold}, confirming that the trained selector loads
and ranks inside the serving path, not only in an offline harness.

\paragraph{Content-addressed cache reuse.} The content-addressed cache yields a
\textbf{92$\times$ mean speedup} on a reuse micro-benchmark that isolates a warm
cached selection against a cold recomputation. This is a mechanism demonstration:
it shows the categorical warm-versus-cold gap, not an end-to-end task number.

\subsection{Honest negatives --- where the method gives nothing}
The boundary of the method is as important as its wins, and we report it as
measured failure, not omission. Table~\ref{tab:neg} lists the negatives.

\begin{table}[h]
\centering
\begin{tabular}{lll}
\toprule
Task & Result & Reading \\
\midrule
Rerank a language model's generations & lift $+0.000$ & = random \\
Public 101-route retrieval set & keyword wins & baseline beats energy \\
\bottomrule
\end{tabular}
\caption{Reproduced negatives. The method adds no signal to free-form generation,
and is not universally better than a lexical baseline on retrieval.}
\label{tab:neg}
\end{table}

\paragraph{Free-form generation: nothing.} Using the energy head to rerank a
free-form language model's own candidate generations produces a \textbf{lift of
$+0.000$ --- indistinguishable from random selection}. The energy carries no signal
about free-form answer quality. This is the sharp edge of the method: it ranks
\emph{discrete candidates with separable structure}, and a model's prose
continuations are neither.

\paragraph{Public retrieval: a keyword baseline can win.} On a public 101-route
retrieval set, a \textbf{plain keyword baseline beats the energy ranker}. We do not
hide this behind the favourable route-ranking number above. The energy ranker's
advantage is not universal; on some retrieval distributions a lexical match is
simply better, and the method should not be deployed as if it dominates everywhere.

\paragraph{The boundary, stated plainly.} Taken together, the wins and the
negatives draw one line: \emph{a learned energy plus a content-addressed cache wins
on discrete structure} --- route retrieval, type classification, cache reuse ---
\emph{and contributes nothing to free-form prose generation}. The selector picks
which route fires; it does not, and should not, replace the language model that
writes text. This boundary is the honest core of the result, and we state it rather
than engineer around it.

\section{Discussion}
The result is deliberately narrow. Energy-based selection is the right tool when the
candidate set is finite, the candidates carry separable structure, and there is an
observed outcome signal to train on --- which is exactly the situation of a runtime
choosing among an agent's learned routes. It is the wrong tool when the ``candidates''
are open-ended text continuations, because the energy has no separable structure to
score; there, the measured lift is zero. Reporting the zero is the point: the
inverse of optimising a metric until it stops meaning anything is naming, precisely,
where the metric goes flat.

Two consequences follow. First, selection and generation are best kept as separate
responsibilities: a discriminative energy head for \emph{which route}, a generative
model for \emph{what to say}. Second, the absolute route-ranking number is a floor,
not a ceiling --- it is a 1-in-100 task with real distractors, and the load-bearing
quantity is the $4.6\times$ multiple over the lexical baseline, which the cache then
makes cheap to serve repeatedly.

\section{Conclusion}
We presented energy-based route ranking: a learned energy head that scores
\emph{(intent, route)} compatibility, ranks an agent's tool library by ascending
energy, and fires the lowest-energy match, trained from observed outcome triples and
embedded into the runtime with a content-addressed cache. On discrete-structure tasks
the method wins and reproduces --- route ranking at $4.6\times$ a keyword baseline,
type classification at $\approx$2.1$\times$, a closed-loop ranker at AUC $0.750$, a
runtime head at $0.83$/$0.64$, and a $92\times$ cache reuse speedup. On free-form
generation it wins nothing (lift $+0.000$), and on at least one public retrieval set
a keyword baseline beats it. The contribution is the line between those two halves:
the method is a discrete-structure selector, and naming where it stops is what makes
the wins trustworthy.

\appendix
\section{Reproducibility}
\label{app:repro}
Every reported number re-runs from a single command. The full set runs with
\texttt{bash bench/energy/reproduce-all.sh}. Individual gates:
\begin{itemize}[leftmargin=1.6em]
  \item Route ranking (R@1, $4.6\times$): \texttt{bash bench/energy/route-ranking-gate.sh}.
  \item Access-pattern type classification ($\approx$2.1$\times$, 2 seeds):
  \texttt{bash bench/energy/intent-type-gate.sh}.
  \item Closed-loop learned ranker (cold-cell AUC $0.750$):
  \texttt{bash bench/ebm-closed-loop-gate.sh}.
  \item Runtime-embedded head (warm $0.83$ / cold $0.64$):
  \texttt{bash bench/ebm-runtime-ship-gate.sh}.
  \item Content-addressed cache reuse ($92\times$):
  \texttt{python3 paper/reference/bench/bench\_reuse.py}.
\end{itemize}
Each gate exits $0$ only when its measured number clears its threshold, so a
regression fails honestly rather than printing a stale figure. The negatives are
recorded in the same accumulating ledger as the wins, so the boundary is part of the
reproduced record, not an aside.

\begin{thebibliography}{99}
\bibitem{lecun} Y.~LeCun, S.~Chopra, R.~Hadsell, M.~Ranzato, F.~Huang.
\emph{A Tutorial on Energy-Based Learning}. In \emph{Predicting Structured Data},
MIT Press, 2006.
\bibitem{ijepa} M.~Assran et al. \emph{Self-Supervised Learning from Images with a
Joint-Embedding Predictive Architecture}. CVPR, 2023.
\bibitem{dpr} V.~Karpukhin et al. \emph{Dense Passage Retrieval for Open-Domain
Question Answering}. EMNLP, 2020.
\bibitem{ltr} C.~Burges et al. \emph{Learning to Rank using Gradient Descent}.
ICML, 2005.
\bibitem{toolformer} T.~Schick et al. \emph{Toolformer: Language Models Can Teach
Themselves to Use Tools}. NeurIPS, 2023.
\bibitem{gorilla} S.~G.~Patil et al. \emph{Gorilla: Large Language Model Connected
with Massive APIs}. arXiv:2305.15334, 2023.
\end{thebibliography}

\end{document}
